%% Question bank chapter file - generated by bank_tools.py build (do not edit while the bank is still being built:
%% it is regenerated from the per-test data in Question Bank\data). Every question is one bankq block:
%%   \begin{bankq}[id=..., source={...}, paper=.., qno=.., type=..., subtopic={...}, answer={...}, answer_src=key|solved,
%%                 working={...}, difficulty=1-10, diff_parts={K S I C F}, diff_note={...},
%%                 group=<passage id or empty>, also_in={duplicate copies in other tests}]
%% difficulty = min(10, K+S+I+C+F-1): K concepts 1-3, S steps 1-4, I insight 0-2, C computation 0-2, F format 0-1
%%   ...question LaTeX, options, figures (\bfig / \bfigi -> img/)...
%%   \end{bankq}
%% Passages shared by several questions are bankpassage blocks; their questions carry group=<passage id>.
\documentclass[10pt]{article}
\usepackage[a4paper,margin=1.8cm]{geometry}
\usepackage[T1]{fontenc}
\usepackage{amsmath,amssymb,mathtools}
\usepackage{graphicx}
\usepackage[export]{adjustbox}
\usepackage{xcolor,needspace,keyval}
\usepackage{enumitem}
\usepackage{booktabs,array,tabularx}
\usepackage[hidelinks]{hyperref}
\setlength{\parindent}{0pt}
\setlength{\parskip}{3pt}
\makeatletter
\define@key{bq}{id}{\def\bq@id{#1}}\define@key{bq}{source}{\def\bq@src{#1}}\define@key{bq}{paper}{\def\bq@paper{#1}}
\define@key{bq}{qno}{\def\bq@qno{#1}}\define@key{bq}{type}{\def\bq@type{#1}}\define@key{bq}{subtopic}{\def\bq@sub{#1}}
\define@key{bq}{answer}{\def\bq@ans{#1}}\define@key{bq}{answer_src}{\def\bq@asrc{#1}}\define@key{bq}{working}{\def\bq@work{#1}}
\define@key{bq}{group}{\def\bq@group{#1}}\define@key{bq}{also_in}{\def\bq@also{#1}}
\define@key{bq}{difficulty}{\def\bq@diff{#1}}\define@key{bq}{diff_parts}{\def\bq@dparts{#1}}\define@key{bq}{diff_note}{\def\bq@dnote{#1}}
\newenvironment{bankq}[1][]{\def\bq@work{}\def\bq@also{}\def\bq@group{}\def\bq@diff{?}\def\bq@dparts{}\def\bq@dnote{}\setkeys{bq}{#1}%
  \par\needspace{4\baselineskip}\medskip\noindent{\small\textbf{\bq@id}\hfill\color{gray}[\bq@src, Paper \bq@paper, Q\bq@qno]
  \bq@type\ $\cdot$ \bq@sub\ $\cdot$ \textbf{D\bq@diff}}\par\nopagebreak\relax}
  {\par{\small\color{teal}Answer: \textbf{\bq@ans}\ifx\bq@work\@empty\else\ (solved: \bq@work)\fi
  \ifx\bq@also\@empty\else\par\color{gray}Also in: \bq@also\fi}\par}
\define@key{bp}{id}{\def\bp@id{#1}}\define@key{bp}{source}{\def\bp@src{#1}}\define@key{bp}{questions}{\def\bp@qs{#1}}
\newenvironment{bankpassage}[1][]{\setkeys{bp}{#1}\par\needspace{4\baselineskip}\bigskip\noindent
  {\small\textbf{Passage \bp@id}\hfill\color{gray}[\bp@src; questions \bp@qs]}\par\nopagebreak\relax\itshape}{\par}
\makeatother
\newcommand{\bfig}[2]{\par\smallskip{\centering\includegraphics[width=#1pt,max width=0.62\linewidth,max height=0.32\textheight]{img/#2}\par}\smallskip}
\newcommand{\bfigi}[2]{\includegraphics[width=#1pt,max width=\linewidth,max height=0.32\textheight,valign=c]{img/#2}}
\newcommand{\bmissing}[1]{\textit{[figure missing in the source conversion: \detokenize{#1}]}}
\newenvironment{qoptions}{\begin{itemize}[leftmargin=1.8em,itemsep=1pt,topsep=2pt,label={}]}{\end{itemize}}
\newcommand{\opt}[2]{\item[(#1)] #2}
\providecommand{\fwstart}{}\providecommand{\fwend}{}\providecommand{\qfullsplit}{}
\newcolumntype{L}[1]{>{\raggedright\arraybackslash}p{#1}}
\newcolumntype{C}[1]{>{\centering\arraybackslash}p{#1}}
\newcolumntype{R}[1]{>{\raggedleft\arraybackslash}p{#1}}
\providecommand{\sectionhdr}[1]{}\providecommand{\subsectionhdr}[1]{}


\begin{document}
{\Large\bfseries Redox Reactions --- Question Bank}\par\smallskip
{\small Chemistry chapter C08. Questions from all converted test papers, grouped by source series. Answers: printed key where it matches the paper, otherwise solved (answer\_src=solved, with working). Exact repeats across papers appear once, with also\_in listing the other copies.}\par\medskip\hrule

\section*{TOAS papers}
\begin{bankq}[id=TOAS03-P1-CHEM-S3-Q03, source={TOAS - Test 03 (27 Feb 2022)}, paper=1, qno=3, type=numerical, subtopic={Equivalents}, answer={135}, answer_src=key, difficulty=5, diff_parts={K2 S2 I1 C0 F1}, diff_note={concepts: non-stoichiometric oxide, redox titration | steps: Fe2+ and Fe3+ composition of Fe0.8O; millimoles of permanganate and Fe2+ oxidised; solve for n; Fe3+ total | I: charge balance in Fe0.8O}]
A 1 litre solution containing equal moles of FeO and $\mathrm{Fe_{0.8}O}$ was titrated with 70 ml 0.3M $\mathrm{KMnO_4}$ in acidic medium. Millimoles of $\mathrm{Fe^{3+}}$ produced are -
\end{bankq}
\begin{bankq}[id=TOAS05-P1-CHEM-S3-Q05, source={TOAS - Test 05 (03 Apr 2022)}, paper=1, qno=5, type=numerical, subtopic={Equivalents}, answer={240.00}, answer_src=key, difficulty=4, diff_parts={K1 S2 I1 C0 F1}, diff_note={concepts: permanganate in acidic and neutral medium | steps: n-factor 5 in acid and 3 in neutral; mmol of MnO4- for the oxalate part; mmol for the iodide part; add volumes | I: neutral medium gives MnO2}]
An aq. solution of 0.5M $\mathrm{KMnO_4}$ is divided into two parts. One part of it requires 125 ml of 1.5M aq. solution of oxalate ions in acidic medium, while another part requires 270 ml of 0.5M aq. solution of iodide ions in neutral medium which are converted into $\mathrm{I_2}$ only. Calcualte total volume (mL) of the initial $\mathrm{KMnO_4}$ solution.
\end{bankq}
\begin{bankq}[id=TOAS05-P2-CHEM-S1-Q13, source={TOAS - Test 05 (03 Apr 2022)}, paper=2, qno=13, type=single, subtopic={Balancing}, answer={C}, answer_src=key, difficulty=1, diff_parts={K1 S1 I0 C0 F0}, diff_note={concepts: balancing a redox equation | steps: balance Cu with HNO3; count N atoms and water}]
The number of Nitrogen atoms, water molecules and the total charge on the product side are respectively when following reaction is balanced ?\\
$\mathrm{Cu+HNO_3\to Cu(NO_3)_2+NO_2+H_2O}$
\begin{qoptions}
\opt{A}{6, 3, 0}
\opt{B}{4, 2, 2}
\opt{C}{4, 2, 0}
\opt{D}{3, 2, 0}
\end{qoptions}
\end{bankq}
\begin{bankq}[id=TOAS06-P1-CHEM-S2-Q05, source={TOAS - Test 06 (17 Apr 2022)}, paper=1, qno=5, type=multi, subtopic={Equivalents}, answer={B,D}, answer_src=key, difficulty=4, diff_parts={K1 S2 I1 C0 F1}, diff_note={concepts: back titration with KIO3 | steps: moles of KI per mL from the standardisation; KI left in the filtrate; KI consumed by AgNO3 gives its percentage | I: KIO3 oxidises iodide to ICl}]
One gram of commercial $\mathrm{AgNO_3}$ sample is dissolved in 50 ml of water. If is treated with 50 ml of a KI solution. The silver iodide thus precipitated is filtered off. Excess of KI in the filtrate is titrated with (M/10) $\mathrm{KIO_3}$ solution in presence of 6 M HCl till all $\mathrm{I^-}$ ions are converted into ICl. It requires 50 mL of (M/10) $\mathrm{KIO_3}$ solution, 20 mL of the same stock solution of KI requires 30 mL of (M/10) $\mathrm{KIO_3}$ under similar conditions.
\begin{qoptions}
\opt{A}{The \% of $\mathrm{AgNO_3}$ in the sample $=50\%$}
\opt{B}{The \% of $\mathrm{AgNO_3}$ in the sample $=85\%$}
\opt{C}{Mass of $\mathrm{AgNO_3}$ in the sample $=0.5$ g}
\opt{D}{Mass of $\mathrm{AgNO_3}$ in the sample $=0.85$ g}
\end{qoptions}
\end{bankq}
\begin{bankq}[id=TOAS08-P1-CHEM-S1-Q03, source={TOAS - Test 08 (15 May 2022)}, paper=1, qno=3, type=multi, subtopic={Equivalents}, answer={A,B,D}, answer_src=key, difficulty=3, diff_parts={K1 S2 I0 C0 F1}, diff_note={concepts: equivalents in redox reactions | steps: n-factor of Cu2S, CuS and S2O3 2-; moles of iodine; volume of each solution}]
250 ml of 1M $\mathrm{I_2}$ is separately reacted with 0.5M $\mathrm{Cu_2S}$ solution, 0.5M $\mathrm{CuS}$ solution and 0.5 M $\mathrm{S_2O_3^{2-}}$ solution causing production of $\mathrm{Cu^{+2}}$ and $\mathrm{SO_4^{2-}}$ in the first two and $\mathrm{S_4O_6^{2-}}$ in the last case along with iodide ions. Which of the following option is/are correct assuming 100\% completion of the reaction. ($\mathrm{I}=127$)
\begin{qoptions}
\opt{A}{100 ml of $\mathrm{Cu_2S}$ solution will be consumed.}
\opt{B}{125 ml of $\mathrm{CuS}$ solution will be consumed.}
\opt{C}{500 ml of $\mathrm{S_2O_3^{2-}}$ solution will be consumed.}
\opt{D}{Equivalent weight of $\mathrm{I_2}$ in each of the reaction will be 127.}
\end{qoptions}
\end{bankq}
\end{document}
